\documentclass[11pt,a4paper]{article}

\usepackage[utf8]{inputenc}
\usepackage[T1]{fontenc}
\usepackage[spanish,es-noshorthands]{babel}
\usepackage{xcolor}
\usepackage[most]{tcolorbox}
\usepackage{listings}
\usepackage{fancyhdr}
\usepackage{lastpage}
\usepackage{enumitem}
\usepackage[hidelinks]{hyperref}

\setlength{\headheight}{30pt}
\pagestyle{fancy}
\fancyhf{}
\renewcommand{\headrulewidth}{0pt}
\fancyhead[R]{%
  \fontsize{8}{9.6}\selectfont
  Licenciatura en Ciencias de la Computación\\
  Sistemas Operativos\\[2pt]
  \rule{4.2cm}{0.5pt}
}
\fancyfoot[C]{\thepage\ / \pageref{LastPage}}

\newtcolorbox{trackbox}[1]{
  colback=green!7!white, colframe=green!45!black,
  arc=2mm, boxrule=0.6pt, left=8pt, right=8pt, top=6pt, bottom=6pt,
  title=TRACK -- #1, fonttitle=\bfseries, coltitle=green!35!black,
  colbacktitle=green!7!white, boxsep=1pt, breakable
}

\newtcolorbox{extrabox}{
  colback=orange!8!white, colframe=orange!55!white,
  arc=2mm, boxrule=0.6pt, left=8pt, right=8pt, top=6pt, bottom=6pt,
  title=EXTRA, fonttitle=\bfseries, coltitle=orange!55!black,
  colbacktitle=orange!8!white, boxsep=1pt, breakable
}

\lstdefinestyle{codestyle}{
  backgroundcolor=\color{black!5},
  rulecolor=\color{black!25},
  frame=single,
  framerule=0.5pt,
  basicstyle=\ttfamily\small,
  breaklines=true,
  showstringspaces=false,
  tabsize=4,
  columns=fullflexible,
}
\lstset{style=codestyle}

\title{Trabajo Práctico 1 -- Ejercicios extra}
\author{}
\date{}

\begin{document}
\renewcommand{\contentsname}{Tracks}
\maketitle
\thispagestyle{fancy}
\begin{center}
SISTEMAS OPERATIVOS
\end{center}
\newpage

\tableofcontents
\newpage

\section{Archivos y llamadas al sistema (3/10)}

\begin{trackbox}{3/10}
\textbf{Descripción:} cómo los programas acceden a archivos a través de llamadas al sistema. Se utilizan C, Python y código fuente de comandos de Linux.
\end{trackbox}

\begin{enumerate}
  \item Leer un archivo en C usando \texttt{open} y \texttt{read}.
  \item Repetir el ejercicio en Python usando \texttt{open} y \texttt{read}.
  \item Copiar un archivo en C usando \texttt{open}, \texttt{read} y \texttt{write}.
  \item Implementar la misma copia en Python.
  \item Encontrar el código fuente de \texttt{ls} y \texttt{cat}.
  \item Identificar cómo abren, leen y muestran archivos.
\end{enumerate}

\newpage
\section{Pipelines (7/10)}

\begin{trackbox}{7/10}
\textbf{Descripción:} cómo conectar comandos y transformar datos usando streams. Se utilizan pipes, redirecciones, \texttt{grep}, \texttt{sort}, \texttt{uniq}, \texttt{wc}, \texttt{xargs}, \texttt{sed}, \texttt{awk}, \texttt{curl}, \texttt{wget} y \texttt{jq}; también se exploran sus opciones mediante las páginas de manual.
\end{trackbox}

\begin{enumerate}
  \item Redirigir la salida con \texttt{>} y \texttt{>>}.
  \item Conectar comandos usando \texttt{|}.
  \item Filtrar resultados con \texttt{grep}.
  \item Ordenar, eliminar duplicados y contar con \texttt{sort}, \texttt{uniq} y \texttt{wc}.
  \item Ejecutar un comando por cada resultado usando \texttt{xargs}.
  \item Editar texto en un stream usando \texttt{sed}.
  \item Extraer campos y calcular valores usando \texttt{awk}.
  \item Descargar datos de premios Nobel:
    \begin{lstlisting}[language=bash]
curl -o prizes.json https://api.nobelprize.org/v1/prize.json
    \end{lstlisting}
  \item Procesar el JSON con \texttt{jq}:
    \begin{itemize}
      \item contar premios por año y categoría;
      \item listar premiados de un año;
      \item contar nombres que comienzan con una letra;
      \item encontrar premiados múltiples.
    \end{itemize}
\end{enumerate}

\newpage
\section{Inspección del sistema (5/10)}

\begin{trackbox}{5/10}
\textbf{Descripción:} cómo Linux expone el estado del sistema. Se utilizan \texttt{/proc}, \texttt{ps}, \texttt{top}, \texttt{free}, \texttt{df} y \texttt{du}; también se exploran sus opciones mediante las páginas de manual.
\end{trackbox}

\begin{enumerate}
  \item Inspeccionar información de CPU en \texttt{/proc/cpuinfo}.
  \item Inspeccionar memoria en \texttt{/proc/meminfo}.
  \item Inspeccionar uptime y load average.
  \item Inspeccionar filesystems montados y uso de disco.
  \item Explorar \texttt{/proc/<pid>} para procesos activos.
  \item Usar \texttt{ps}, \texttt{top}, \texttt{free}, \texttt{df} y \texttt{du}.
  \item Comparar la salida de estos comandos con la información de \texttt{/proc}.
  \item Escribir un script que extraiga valores relevantes y muestre un reporte conciso. Este ejercicio conecta con el track de Pipelines.
\end{enumerate}

\newpage
\section{Entorno del shell (5/10)}

\begin{trackbox}{5/10}
\textbf{Descripción:} cómo el shell configura y lanza programas. Se utilizan variables de entorno, \texttt{PATH}, exportación, herencia y \texttt{env}; también se exploran sus opciones mediante las páginas de manual.
\end{trackbox}

\begin{enumerate}
  \item Inspeccionar variables usando \texttt{env}, \texttt{printenv} y \texttt{echo}.
  \item Guardar el valor actual de \texttt{PATH}.
  \item Ejecutar \texttt{export PATH=""} y observar qué comandos fallan.
  \item Ejecutar comandos usando rutas absolutas.
  \item Restaurar \texttt{PATH}.
  \item Definir una variable sin exportarla y observar su ausencia en un shell hijo.
  \item Exportar la variable y repetir la prueba.
  \item Modificar \texttt{HOME}, \texttt{PATH} y variables propias usando \texttt{env}.
\end{enumerate}

\newpage
\section{Procesos (3/10)}

\begin{trackbox}{3/10}
\textbf{Descripción:} cómo Linux administra programas en ejecución. Se utilizan jobs, estados de procesos, señales, \texttt{ps}, \texttt{top} y \texttt{/proc}; también se exploran sus opciones mediante las páginas de manual.
\end{trackbox}

\begin{enumerate}
  \item Ejecutar comandos en foreground y background.
  \item Usar \texttt{jobs}, \texttt{fg} y \texttt{bg}.
  \item Inspeccionar procesos con \texttt{ps} y \texttt{top}.
  \item Pausar y reanudar procesos.
  \item Enviar señales usando \texttt{kill}.
  \item Comparar terminación controlada y forzada.
  \item Observar procesos padre e hijo.
  \item Inspeccionar estados en \texttt{/proc/<pid>/status}.
\end{enumerate}

\end{document}
